% GENERATED FILE. Edit canonical lessons, then run npm run generate:books.
% canonical-source-hash: fnv1a64:434538c3b2b7ec3e
% canonical-lessons: SA-C448-bhavah, SA-C448-anutapah, SA-C448-hrsyati, SA-C448-ullasah, SA-C448-khedah

\chapter{Naming a Feeling}
\label{ch:sa-naming-a-feeling}
\hlchaptermodality{448}

\begin{chapteropening}
\textbf{By the end of this chapter:} \emph{I can name a feeling and say that I felt it.}

Complete the last lesson of chapter 448: I can name a feeling and say that I felt it.
\end{chapteropening}

\section[\texorpdfstring{\sk{भावः} (bhāvaḥ)}{भावः (bhāvaḥ)}]{\sk{भावः} (bhāvaḥ) — a feeling; a state of mind}

\label{lesson:SA-C448-bhavah}



\begin{quote}
Before the new one: say the Sanskrit for radio, then the Sanskrit for the internet.
\end{quote}

\subsection*{You'll want to know: \sk{भावः}}
\textbf{\sk{भावः}} — \emph{bhāvaḥ} — \textquotedblleft{}a feeling; a state of mind\textquotedblright{}.

\begin{quote}
\sk{भावः} \sk{हृदये} \sk{जायते।}
\end{quote}

\begin{quote}
\emph{bhāvaḥ hṛdaye jāyate.}
\end{quote}

\begin{quote}
A feeling arises in the heart.
\end{quote}

A masculine noun.

\begin{cousinweb}
\textbf{\sk{भावः}} (\emph{bhāvaḥ}) is built on the root √bhū, \textquotedblleft{}to be, to become\textquotedblright{}; MW gives \emph{bhāva} the senses \textquotedblleft{}a state of mind or body, a way of feeling\textquotedblright{} and, in poetics, \textquotedblleft{}emotion\textquotedblright{} (MW).
\end{cousinweb}

\begin{grammarlens}[title={What you've built}]
One more word for describing what you lived and felt.
\end{grammarlens}

\subsection*{Your turn}
\begin{itemize}
\raggedright
  \item \emph{Say it:} \emph{bhāvaḥ}
  \item \emph{Say it:} \emph{bhāvaḥ}, once more
  \item \emph{Say it:} \emph{bhāvaḥ hṛdaye jāyate.} — A feeling arises in the heart.
  \item \emph{Recall:} say the Sanskrit for radio, then the Sanskrit for the internet, then say \emph{bhāvaḥ} again
\end{itemize}

\begin{tcolorbox}[breakable,colback=teal!4,colframe=teal!35!black,title={Before you move on}]
What does \sk{भावः} mean? (\textquotedblleft{}A feeling; a state of mind\textquotedblright{}.) Say it once more.
\end{tcolorbox}
\section[\texorpdfstring{\sk{अनुतापः} (anutāpaḥ)}{अनुतापः (anutāpaḥ)}]{\sk{अनुतापः} (anutāpaḥ) — regret; repentance}

\label{lesson:SA-C448-anutapah}



\begin{quote}
Before the new one: say the Sanskrit for the internet, then the Sanskrit for a feeling.
\end{quote}

\subsection*{You'll want to know: \sk{अनुतापः}}
\textbf{\sk{अनुतापः}} — \emph{anutāpaḥ} — \textquotedblleft{}regret; repentance\textquotedblright{}.

\begin{quote}
\sk{कलहात्} \sk{अनन्तरं} \sk{मम} \sk{अनुतापः} \sk{भवति।}
\end{quote}

\begin{quote}
\emph{kalahāt anantaraṃ mama anutāpaḥ bhavati.}
\end{quote}

\begin{quote}
After a quarrel I feel regret.
\end{quote}

A masculine noun.

\begin{cousinweb}
\textbf{\sk{अनुतापः}} (\emph{anutāpaḥ}) is built from \emph{anu}, \textquotedblleft{}after\textquotedblright{}, and \emph{tāpa}, \textquotedblleft{}heat\textquotedblright{}, from the root √tap, \textquotedblleft{}to burn\textquotedblright{}, the root of \textbf{\sk{तापः}}: a burning that comes afterwards; MW gives \textquotedblleft{}repentance\textquotedblright{} (MW).
\end{cousinweb}

\begin{grammarlens}[title={What you've built}]
One more word for describing what you lived and felt.
\end{grammarlens}

\subsection*{Your turn}
\begin{itemize}
\raggedright
  \item \emph{Say it:} \emph{anutāpaḥ}
  \item \emph{Say it:} \emph{anutāpaḥ}, once more
  \item \emph{Say it:} \emph{kalahāt anantaraṃ mama anutāpaḥ bhavati.} — After a quarrel I feel regret.
  \item \emph{Recall:} say the Sanskrit for the internet, then the Sanskrit for a feeling, then say \emph{anutāpaḥ} again
\end{itemize}

\begin{tcolorbox}[breakable,colback=teal!4,colframe=teal!35!black,title={Before you move on}]
What does \sk{अनुतापः} mean? (\textquotedblleft{}Regret; repentance\textquotedblright{}.) Say it once more.
\end{tcolorbox}
\section[\texorpdfstring{\sk{हृष्यति} (hṛṣyati)}{हृष्यति (hṛṣyati)}]{\sk{हृष्यति} (hṛṣyati) — thrills with joy}

\label{lesson:SA-C448-hrsyati}



\begin{quote}
Before the new one: say the Sanskrit for a feeling, then the Sanskrit for regret.
\end{quote}

\subsection*{You'll want to know: \sk{हृष्यति}}
\textbf{\sk{हृष्यति}} — \emph{hṛṣyati} — \textquotedblleft{}thrills with joy\textquotedblright{}.

\begin{quote}
\sk{बालः} \sk{उपहारं} \sk{पश्यति}, \sk{हृष्यति} \sk{च।}
\end{quote}

\begin{quote}
\emph{bālaḥ upahāraṃ paśyati, hṛṣyati ca.}
\end{quote}

\begin{quote}
The child sees the present and thrills with joy.
\end{quote}

Past (he …): \sk{सः} \sk{हृष्टवान्}, \emph{saḥ hṛṣṭavān}.

\begin{cousinweb}
\textbf{\sk{हृष्यति}} (\emph{hṛṣyati}) comes from the root √hṛṣ, \textquotedblleft{}to bristle, to thrill\textquotedblright{}, the root of \textbf{\sk{हर्षः}}: delight felt as the hair standing on end; MW compares Latin \emph{horreo} (MW).
\end{cousinweb}

\begin{grammarlens}[title={What you've built}]
One more word for describing what you lived and felt.
\end{grammarlens}

\subsection*{Your turn}
\begin{itemize}
\raggedright
  \item \emph{Say it:} \emph{hṛṣyati}
  \item \emph{Say it:} \emph{hṛṣyati}, once more
  \item \emph{Say it:} \emph{bālaḥ upahāraṃ paśyati, hṛṣyati ca.} — The child sees the present and thrills with joy.
  \item \emph{Recall:} say the Sanskrit for a feeling, then the Sanskrit for regret, then say \emph{hṛṣyati} again
\end{itemize}

\begin{tcolorbox}[breakable,colback=teal!4,colframe=teal!35!black,title={Before you move on}]
What does \sk{हृष्यति} mean? (\textquotedblleft{}Thrills with joy\textquotedblright{}.) Say it once more.
\end{tcolorbox}
\section[\texorpdfstring{\sk{उल्लासः} (ullāsaḥ)}{उल्लासः (ullāsaḥ)}]{\sk{उल्लासः} (ullāsaḥ) — merriment; high spirits}

\label{lesson:SA-C448-ullasah}



\begin{quote}
Before the new one: say the Sanskrit for regret, then the Sanskrit for thrills with joy.
\end{quote}

\subsection*{You'll want to know: \sk{उल्लासः}}
\textbf{\sk{उल्लासः}} — \emph{ullāsaḥ} — \textquotedblleft{}merriment; high spirits\textquotedblright{}.

\begin{quote}
\sk{उत्सवे} \sk{सर्वत्र} \sk{उल्लासः} \sk{अस्ति।}
\end{quote}

\begin{quote}
\emph{utsave sarvatra ullāsaḥ asti.}
\end{quote}

\begin{quote}
At the festival there is merriment everywhere.
\end{quote}

A masculine noun.

\begin{cousinweb}
\textbf{\sk{उल्लासः}} (\emph{ullāsaḥ}) is built from \emph{ud}, \textquotedblleft{}up\textquotedblright{}, and the root √las, \textquotedblleft{}to shine, to play\textquotedblright{}; MW gives \textquotedblleft{}light\textquotedblright{} first and then \textquotedblleft{}joy, merriness\textquotedblright{} (MW).
\end{cousinweb}

\begin{grammarlens}[title={What you've built}]
One more word for describing what you lived and felt.
\end{grammarlens}

\subsection*{Your turn}
\begin{itemize}
\raggedright
  \item \emph{Say it:} \emph{ullāsaḥ}
  \item \emph{Say it:} \emph{ullāsaḥ}, once more
  \item \emph{Say it:} \emph{utsave sarvatra ullāsaḥ asti.} — At the festival there is merriment everywhere.
  \item \emph{Recall:} say the Sanskrit for regret, then the Sanskrit for thrills with joy, then say \emph{ullāsaḥ} again
\end{itemize}

\begin{tcolorbox}[breakable,colback=teal!4,colframe=teal!35!black,title={Before you move on}]
What does \sk{उल्लासः} mean? (\textquotedblleft{}Merriment; high spirits\textquotedblright{}.) Say it once more.
\end{tcolorbox}
\section[\texorpdfstring{\sk{खेदः} (khedaḥ)}{खेदः (khedaḥ)}]{\sk{खेदः} (khedaḥ) — distress; weariness of heart}

\label{lesson:SA-C448-khedah}



\begin{quote}
Before the new one: say the Sanskrit for thrills with joy, then the Sanskrit for merriment.
\end{quote}

\subsection*{You'll want to know: \sk{खेदः}}
\textbf{\sk{खेदः}} — \emph{khedaḥ} — \textquotedblleft{}distress; weariness of heart\textquotedblright{}.

\begin{quote}
\sk{मम} \sk{मित्रं} \sk{रुग्णम्}, \sk{अतः} \sk{मम} \sk{खेदः} \sk{अस्ति।}
\end{quote}

\begin{quote}
\emph{mama mitraṃ rugṇam, ataḥ mama khedaḥ asti.}
\end{quote}

\begin{quote}
My friend is ill, so I am distressed.
\end{quote}

A masculine noun.

\begin{cousinweb}
\textbf{\sk{खेदः}} (\emph{khedaḥ}) comes from the root √khid, \textquotedblleft{}to be depressed\textquotedblright{}; MW gives \textquotedblleft{}lassitude, depression\textquotedblright{} and \textquotedblleft{}pain, distress\textquotedblright{} (MW).
\end{cousinweb}

\begin{grammarlens}[title={What you've built}]
One more word for describing what you lived and felt.
\end{grammarlens}

\subsection*{Your turn}
\begin{itemize}
\raggedright
  \item \emph{Say it:} \emph{khedaḥ}
  \item \emph{Say it:} \emph{khedaḥ}, once more
  \item \emph{Say it:} \emph{mama mitraṃ rugṇam, ataḥ mama khedaḥ asti.} — My friend is ill, so I am distressed.
  \item \emph{Recall:} say the Sanskrit for thrills with joy, then the Sanskrit for merriment, then say \emph{khedaḥ} again
\end{itemize}

\begin{tcolorbox}[breakable,colback=teal!4,colframe=teal!35!black,title={Before you move on}]
What does \sk{खेदः} mean? (\textquotedblleft{}Distress; weariness of heart\textquotedblright{}.) Say it once more.
\end{tcolorbox}
